\chapter{اللقاء السادس عشر: تابع تفسير سورة البقرة (٥)}

\section{مقدمة وتصحيح}
السلام عليكم ورحمة الله وبركاته. الحمد لله رب العالمين، وأشهد أن لا إله إلا الله وحده لا شريك له، وأشهد أن محمداً عبده ورسوله. اللهم صل على محمد وعلى آل محمد كما صليت على آل إبراهيم إنك حميد مجيد، اللهم بارك على محمد وعلى آل محمد كما باركت على آل إبراهيم إنك حميد مجيد.

هذا هو الدرس الرابع عشر (السادس عشر بترقيم الدورة) من اجتماعنا على دراسة كتاب تفسير \rawi{ابن جرير الطبري} عليه رحمة الله. 

أحب أن أنبه قبل أن نبدأ أن قراءتنا لهذا الكتاب -وطبعاً هذا الكتاب كبير يعني تقريباً ستة وعشرون مجلداً- أنا أعلق تعليقات خفيفة في المواضع التي أظن أنها بحاجة إلى بيان، ولكن بطل هذه القصة هو أنت. يجب أن تقرأ الدرس مرة أخرى، وأن تستمع إليه، وأن تتدارسه مع أصحابك، فهذه هي البداية مع الكتاب وليست النهاية. أحب منك إذا دخلت على كتاب لتدرسه أن تكرر قراءة الكتاب أكثر من مرة، خصوصاً الكتب التأسيسية في بابها، وأن تتدارس هذا الكتاب وتستمع إلى شروحه. 

ولا تعتمد فقط على مجرد هذا الدرس، فنحن أحياناً نفوت بعض المواضع لا نعلق عليها، وأحياناً نعطيكم واجبات حتى تتدربوا على القراءة والتلخيص، وأحياناً نمرر بعض الآثار لا نقرأها اختصاراً، فإذا نحن فقط فتحنا الكتاب، ولكن لا بد أن يكون لك مع الكتاب وقفة أخرى. بارك الله فيكم.

\begin{fawaid}[تصحيح نسبة قول في الإعراب]
في الدرس الماضي، عند تفسير قوله تعالى: \begin{ayah}أَوَكُلَّمَا عَاهَدُوا عَهْدًا\end{ayah}، نُسب القول بأن الواو حرف عطف أُدخلت عليها ألف الاستفهام إلى \rawi{سيبويه} على أنه من نحويي الكوفة. والصواب أن النسبة إلى \rawi{سيبويه} صحيحة، ولكنه من مؤسسي مدرسة البصرة وليس من الكوفيين (وأقرب أهل الكوفة إلى هذا الكلام هو \rawi{الزجاج}).
\end{fawaid}

\separator

\section{تأويل (وما يعلمان من أحد حتى يقولا إنما نحن فتنة فلا تكفر)}
وصلنا بحمد الله تبارك وتعالى في المجلد الثاني إلى القول في تأويل قوله جل ثناؤه: \begin{ayah}وَمَا يُعَلِّمَانِ مِنْ أَحَدٍ حَتَّىٰ يَقُولَا إِنَّمَا نَحْنُ فِتْنَةٌ فَلَا تَكْفُرْ\end{ayah}.

\isnad{قال الطبري رحمه الله:} «وتأويل ذلك وما يعلم الملكان من أحد من الناس الذي أنزل عليهما من التفريق بين المرء وزوجه حتى يقولا له إنما نحن بلاء وفتنة لبني آدم فلا تكفر».

\isnad{كما حدثني} \rawi{موسى} \isnad{قال بإسناده عن} \rawi{السدي} \isnad{قال:} «إذا أتاهما -يعني هاروت وماروت- إنسان يريد السحر وعظاه وقالا له: لا تكفر إنما نحن فتنة. فإذا أبى قالا له: ائت هذا الرماد فبل عليه. فإذا بال عليه خرج منه نور يسطع حتى يدخل السماء وذلك الإيمان. وأقبل شيء أسود كهيئة الدخان حتى يدخل في مسامعه وكل شيء منه فذلك غضب الله، فإذا أخبرهما بذلك علماه السحر، فذلك قول الله: \begin{ayah}وَمَا يُعَلِّمَانِ مِنْ أَحَدٍ حَتَّىٰ يَقُولَا إِنَّمَا نَحْنُ فِتْنَةٌ فَلَا تَكْفُرْ\end{ayah}».

\isnad{وأيضاً قال بإسناده عن} \rawi{قتادة} \rawi{والحسن} \isnad{قالا:} «أخذ عليهما ألا يعلما أحداً حتى يقولا إنما نحن فتنة فلا تكفر».

\subsection{تلخيص أقوال المفسرين في (وما أنزل على الملكين)}
كان من الواجب في الدرس الماضي تلخيص قول الإمام \rawi{الطبري} في معنى (ما) في الآية: \begin{ayah}وَمَا أُنزِلَ عَلَى الْمَلَكَيْنِ\end{ayah}. وفيها أربعة أقوال:

\textbf{القول الأول:} أن (ما) بمعنى الجحد (أي: لم ينزل)، وهو قول \rawi{ابن عباس} و\rawi{الربيع بن أنس}. ويكون التأويل: اتبعوا ما تتلو الشياطين على ملك سليمان، وما كفر سليمان، ولا أنزل الله السحر على الملكين (جبريل وميكائيل). 
وقد بين الإمام \rawi{الطبري} وجه خطأ هذا القول: أنه لو كان المقصود هاروت وماروت أنهما رجلان من بني آدم، لكان يجب أن يرتفع السحر بهلاكهما، وهذا لم يحصل.

\textbf{القول الثاني:} أن (ما) بمعنى (الذي)، وهو قول \rawi{قتادة} و\rawi{الزهري} و\rawi{السدي} و\rawi{ابن زيد}. ويكون التأويل: اتبعت اليهود الذي تلت الشياطين في ملك سليمان والذي أنزل على الملكين. 
وقد دفعوا الإشكال في جواز تعليم الملائكة السحر بأن الله أنزل الخير والشر ليمتحن به العباد، وأنه ليس في مجرد العلم بالسحر إثم، وأن الملكين يخبران من يأتيهما بأنهما فتنة وينهيانه عن السحر. وهذا القول هو الذي رجحه الإمام \rawi{الطبري}.

\textbf{القول الثالث:} أن (ما) بمعنى (الذي) ولكنها عطف على (ما) الأولى، والمعنى: وما أنزل على الملكين أي التفريق بين المرء وزوجه. قال به \rawi{مجاهد}.

\textbf{القول الرابع:} جواز أن تكون (ما) بمعنى (الذي) أو بمعنى الجحد، قاله \rawi{القاسم بن محمد}.

\begin{fawaid}[رد الطبري على جعل (ما) للجحد]
بين الإمام \rawi{الطبري} خطأ توجيه (ما) بمعنى الجحد، موضحاً أنه إن وُجهت لمعنى الجحد (لم ينزل عليهما)، فلا يخلو اسما (هاروت وماروت) من أن يكونا بدلاً من (الملكين) أو بدلاً من (الناس). 
فإن جعلا بدلاً من الملكين، بطل معنى قوله: \begin{ayah}وَمَا يُعَلِّمَانِ مِنْ أَحَدٍ\end{ayah}، فإذا كان الملكان لم ينزل عليهما شيء فبما يفرقون بين المرء وزوجه؟ وما الذي يتعلم منهما؟
وإن جعلا بدلاً من الناس، وجب أن تكون الشياطين هي التي تعلم هاروت وماروت السحر. فلو كانا ملكين، لكانت نسبتهما إلى تعلم السحر من الشياطين وتعليمه الناس أعظم معصية، وهذا يتناقض مع قولهما: \begin{ayah}إِنَّمَا نَحْنُ فِتْنَةٌ فَلَا تَكْفُرْ\end{ayah}. ولو كانا رجلين من بني آدم، لوجب أن يرتفع السحر من العالم بهلاكهما.
\end{fawaid}

\begin{fawaid}[طريقة دراسة كتب الأئمة المحققين]
هذه الطريقة في الدراسة هي أفضل طريقة يمكن أن يُدرس بها كتاب من كتب المحققين. لأن الطالب قد يتوه مع كثرة الأقوال والأدلة والردود. فالقراءة المجردة لا تكفي، بل لا بد من:
\begin{itemize}
    \item \textbf{التصنيف:} أن يعرف الطالب موضع تصوير المسألة، وموضع إيراد الأقوال، وموضع ذكر حجة كل قول، وموضع الردود، وموضع الترجيح.
    \item \textbf{الخلاصة:} أن يدرك خلاصة قول المؤلف بحجته، وكيف أجاب على الإشكالات وكيف رد على المخالفين.
\end{itemize}
وهذا هو الفرق الدقيق بين القارئ العادي وطالب العلم الذي يُرجى له الإتقان. فمن يقرأ ليفهم فقط قد يغيب عنه تصنيف الأقوال واستخراج الخلاصات، بينما الناقد يختبر ويميز الأقوال، ولا يُشترط أن يقطع بالنتائج فوراً، بل يكون ترجيحه ثمرة خبرة ومعالجة للمسائل عبر السنين.
\end{fawaid}

\separator

\section{تأويل (فيتعلمون منهما ما يفرقون به بين المرء وزوجه)}
\isnad{قال أبو جعفر:} «وقوله جل ثناؤه \textit{فيتعلمون منهما} خبر مبتدأ عن المتعلمين من الملكين ما أنزل عليهما، وليس بجواب لقوله \textit{وما يعلمان من أحد}، بل هو خبر مستأنف فلذلك رُفع. فمعنى الكلام إذاً: وما يعلمان من أحد حتى يقولا إنما نحن فتنة، فيأبون قبول ذلك منهما، فيتعلمون منهما ما يفرقون به بين المرء وزوجه».

\isnad{وقال:} «وقد قيل إن قوله \textit{فيتعلمون} خبر عن اليهود معطوف على قوله \textit{ولكن الشياطين كفروا يعلمون الناس السحر وما أنزل على الملكين}... وجعلوا ذلك من المؤخر الذي معناه التقديم. والذي قلنا أشبه بتأويل الآية؛ لأن إلحاق الكلام بالذي يليه من الكلام ما كان للتأويل وجه صحيح أولى من إلحاقه بما قد حيل بينه وبينه من معترض الكلام».

\begin{fawaid}[قاعدة في الترجيح: إلحاق الكلام بما يليه أولى]
هذه قاعدة عند \rawi{الطبري} في الترجيح بين الأقوال: إلحاق الكلام بالذي يليه (يتصل به مباشرة) أولى من إلحاقه بما تباعد عنه وحيل بينهما بكلام معترض، طالما كان للتأويل المتصل وجه صحيح.
\end{fawaid}

\isnad{قال الطبري:} «والهاء والميم والألف من قوله \textit{منهما} من ذكر الملكين. ومعنى ذلك: فيتعلم الناس من الملكين السحر الذي يفرقون به بين المرء وزوجه... وأما (المرء) فإنه بمعنى رجل من أسماء بني آدم... وأما (الزوج) فإن أهل الحجاز يقولون لامرأة الرجل هي زوجه... وتميم وكثير من قيس وأهل نجد يقولون هي زوجته».

\begin{fawaid}[إطلاق لفظ الزوج على المرأة لغةً]
الأفصح في لغة العرب (وهي لغة أهل الحجاز، وبها نزل القرآن) إطلاق لفظ (الزَّوْج) على الذكر والأنثى معاً بلا تاء تأنيث، كما قال تعالى: \textit{أمسك عليك زوجك}. أما إثبات التاء (زوجة) فهي لغة قبائل تميم وقيس.
\end{fawaid}

\subsection{كيف يفرق الساحر بين المرء وزوجه؟}
\isnad{قال الطبري:} «فإن قال قائل: وكيف يفرق الساحر بين المرء وزوجه؟ قد دللنا فيما مضى على أن معنى السحر تخييل الشيء إلى المرء بخلاف ما هو به في عينه وحقيقته... فتفريقه بين المرء وزوجه تخييله بسحره إلى كل واحد منهما شخص الآخر على خلاف ما هو به في حقيقته من حسن وجمال حتى يقبّحه عنده فينصرف بوجهه ويعرض عنه، حتى يُحدث الزوج لامرأته فراقاً. فيكون الساحر مفرقاً بينهما بإحداثه السبب الذي كان عنه فرقة ما بينهما. وقد دللنا... على أن العرب تضيف الشيء إلى مسببه وإن لم يكن باشر فعل ما حدث عن السبب».

\begin{fawaid}[قدرة الساحر وحقيقة السحر]
نلاحظ أن \rawi{الطبري} رحمه الله صرح هنا بأن التفريق بين المرء وزوجه يكون بـ (تخييله) بسحره إلى كل واحد منهما شخص الآخر على خلاف ما هو عليه في حقيقته، حتى يُقبّحه عنده. وكأنه يرى أن السحر هنا ليس تصرفاً في الأعيان أو قدرة مطلقة على تدبير الأمور، بل هو مجرد تخييل يفسد به العلاقة.
\end{fawaid}

\subsection{توجيه الآية على قول المخالفين}
الذين نفوا أن يكون الملكان يُعَلِّمان الناس التفريق بين المرء وزوجه (لأنهم جعلوا "ما" للنفي)، وجهوا تأويل قوله (فيتعلمون منهما) إلى: فيتعلمون \textbf{مكان} ما علماهما، ما يفرقون به بين المرء وزوجه، كأنهم تعلموا الشر مكان الخير.
وهذا من أحسن ما يفعله الطبري أنه يفسر الجملة القرآنية بناءً على كل قول من الأقوال المطروحة.

\separator

\section{تأويل (وما هم بضارين به من أحد إلا بإذن الله)}
\isnad{قال أبو جعفر:} «يعني بقوله جل ثناؤه \textit{وما هم بضارين به من أحد إلا بإذن الله}: وما المتعلمون من الملكين... بضارين بالذي تعلموه منهما... أحداً من الناس إلا من قضى الله عليه أن ذلك يضره. فأما من دفع الله عنه ضره وحفظه من مكروه السحر... فإن ذلك غير ضاره... والإذن في كلام العرب أوجُه: منها الأمر على غير وجه الإلزام... ومنها التخلية بين المأذون له والمخلى بينه وبينه، ومنها العلم بالشيء... كانه قال جل ثناؤه: وما هم بضارين بالذي تعلموا من الملكين من أحد إلا بعلم الله».

\begin{fawaid}[الفرق بين الإذن الكوني والإذن الشرعي]
للإذن والقضاء والأمر والإرادة في القرآن معنيان:
\begin{itemize}
    \item \textbf{معنى شرعي:} وهو ما يحبه الله ويشرعه، مثل: \begin{ayah}أَمْ لَهُمْ شُرَكَاءُ شَرَعُوا لَهُم مِّنَ الدِّينِ مَا لَمْ يَأْذَنْ بِهِ اللَّهُ\end{ayah}.
    \item \textbf{معنى كوني (قدري):} وهو ما يقدره الله بحكمته وإن لم يحبه، مثل: \begin{ayah}وَمَا هُمْ بِضَارِّينَ بِهِ مِنْ أَحَدٍ إِلَّا بِإِذْنِ اللَّهِ\end{ayah}. فهذا الإذن قدري، لأن الله حرّم السحر ولا يحبه، ولكنه أذن بوقوعه لحكمة.
\end{itemize}
وقد قصّر الطبري حين حصر الإذن هنا في (العلم المسبق)، فالأقرب أن الإذن هنا بمعنى (المشيئة والتقدير)، والمشيئة تتضمن العلم والكتابة والخلق، كما قال تعالى: \begin{ayah}وَمَا تَشَاءُونَ إِلَّا أَن يَشَاءَ اللَّهُ\end{ayah}.
\end{fawaid}

\begin{fawaid}[خلق الله للشر نسبي]
الله تبارك وتعالى لا يقدّر شراً محضاً من كل وجه، بل كل مقدر إما أن يكون خيراً محضاً، أو شراً من وجه وخيراً من وجه آخر. كما في حادثة الإفك، كانت شراً على من خاض فيها، ولكنها كانت خيراً في تبرئة عائشة رضي الله عنها وفضح المنافقين وتشريع الأحكام: \begin{ayah}لَا تَحْسَبُوهُ شَرًّا لَّكُم ۖ بَلْ هُوَ خَيْرٌ لَّكُمْ\end{ayah}. ففعل الله كله خير، والشر لا ينسب إليه كما في الحديث: \begin{hadith}والشر ليس إليك\end{hadith}.
\end{fawaid}

\separator

\section{تأويل (ويتعلمون ما يضرهم ولا ينفعهم)}
\isnad{قال أبو جعفر:} «يعني جل ثناؤه بقوله \textit{ويتعلمون}: أي الناس الذين يتعلمون من الملكين... يتعلمون منهما السحر الذي يضرهم في دينهم ولا ينفعهم في معادهم، فأما في العاجل في الدنيا فإنهم قد كانوا يكسبون به ويصيبون به معاشاً».

\begin{fawaid}[خطر العلم الذي لا ينفع]
في هذه الآية معنيان جليلان لطالب العلم:
\begin{itemize}
    \item \textbf{الأول:} ليس كل معرفة نافعة في أصلها، فهناك علوم لا فائدة منها أو هي باب شر كالسحر، والانشغال بها يشغل عن العلم النافع.
    \item \textbf{الثاني:} المعلومات الصحيحة والنافعة قد لا ينتفع بها صاحبها، بل قد تكون سلاحاً يستعمله لنيل أغراضه الدنيوية أو الكبر على الناس أو التبرير للمجرمين، فيكون علمه وبالاً عليه. ولذلك كان النبي ﷺ يستعيذ بالله من علم لا ينفع.
\end{itemize}
\end{fawaid}

\separator

\section{تأويل (ولقد علموا لمن اشتراه ما له في الآخرة من خلاق)}
\isnad{قال أبو جعفر:} «يعني بقوله جل ثناؤه... لقد علم النابذون من يهود بني إسرائيل كتابي وراء ظهورهم تجاهلاً منهم... المؤثرون عليه اتباع السحر الذي تلته الشياطين... أن كل من اشترى السحر بكتابي الذي أنزلته على رسولي، ما له في الآخرة من خلاق».

\begin{fawaid}[استبدال الكفر بالإيمان حسداً]
كان اليهود يعلمون أن النبي ﷺ حق، ولكنهم بدلاً من أن يؤمنوا به ويعزروه، اشتروا الكفر بالإيمان حسداً وبغياً، كما قال تعالى: \begin{ayah}بِئْسَمَا اشْتَرَوْا بِهِ أَنفُسَهُمْ أَن يَكْفُرُوا بِمَا أَنزَلَ اللَّهُ بَغْيًا\end{ayah}. وهذا يعلمنا أن الضلال ليس سببه الجهل دائماً، بل قد يكون الكبر أو الحسد أو حب الدنيا. ومجرد إعطاء النعمة للإنسان (كالعلم) هو ابتلاء، فشكرها يكون باستعمالها فيما يُرضي الله، ومن لم يفعل فقد بدّل نعمة الله كفراً.
\end{fawaid}

\isnad{واختلف أهل التأويل في معنى (الخلاق):} فرجح \rawi{الطبري} أنه "النصيب"، مستدلاً بلغة العرب وبقول النبي ﷺ: \begin{hadith}ليؤيدن الله هذا الدين بأقوام لا خلاق لهم\end{hadith} (أي لا نصيب لهم في الإسلام).

\isnad{وأما إعراب (لَمَنِ اشْتَرَاهُ):} فبيّن \rawi{الطبري} أن (مَنْ) في موضع رفع، لأن قوله (علموا) بمعنى اليمين (والله لَمَن اشترى السحر)، فأجيب بلام اليمين (لَمَنْ).

\begin{fawaid}[إعراب (لمن اشتراه) وتضمين الأفعال]
من لطائف التوجيه النحوي عند الطبري إعراب (مَنْ) في قوله (لَمَنِ اشْتَرَاهُ) في محل رفع مبتدأ، وتوجيه دخول اللام عليها بأن الفعل (علموا) أُشرب معنى اليمين (أي: والله لمن اشتراه)، فجاءت اللام كأنها واقعة في جواب القسم المقدر.
\end{fawaid}

\separator

\section{تأويل (ولبئس ما شروا به أنفسهم لو كانوا يعلمون)}
\isnad{قال أبو جعفر:} «ولبئس ما باع به نفسه من تعلم السحر لو كان يعلم سوء عاقبته».

طرح \rawi{الطبري} إشكالاً: كيف قال (لو كانوا يعلمون) وقد قال قبلها (ولقد علموا)؟
فأجاب بأن ذلك من المؤخر الذي معناه التقديم، أي: وما هم بضارين به من أحد إلا بإذن الله، ويتعلمون ما يضرهم ولا ينفعهم، ولبئس ما شروا به أنفسهم لو كانوا يعلمون، ولقد علموا لمن اشتراه... 
وهذا التوجيه من الطبري لرد دعوى من زعم التناقض، أو من زعم أن "علموا" للشياطين و"يعلمون" للناس (كما زعم الأخفش وردّه الطبري بقوة)، أو من زعم أنهم علموا المعرفة ولكنهم لم يعملوا بها فسُمّوا جاهلين.
وأصّل الطبري قاعدة: \textbf{«تأويل القرآن على المفهوم الظاهر بالخطاب، دون الخفي الباطن منه... حتى تأتي الدلالة من الوجه الذي يجب التسليم له... أولى»}.

\separator

\section{تأويل (ولو أنهم آمنوا واتقوا لمثوبة من عند الله خير)}
\isnad{قال أبو جعفر:} «يعني... لو أن الذين يتعلمون من الملكين... صدقوا الله ورسوله... لكان ثواب الله إياهم على ذلك خيراً لهم من السحر ومما اكتسبوا به لو كانوا يعلمون».

\isnad{والمثوبة:} مصدر من الثواب، وهو ما يرجع إلى العبد بدلاً وعوضاً عن عمله.

\separator

\section{تأويل (يا أيها الذين آمنوا لا تقولوا راعنا وقولوا انظرنا)}
اختلف المفسرون في معنى (راعنا): فقيل: لا تقولوا خلافاً (وهو قول \rawi{مجاهد} و\rawi{عطاء})، وقيل: أَرْعِنا سمعك (أي اسمع منا ونسمع منك)، وقيل: إنها كلمة سخرية باليهودية (رعنا) بمعنى الخطأ والرعونة.

\isnad{قال الطبري مرجحاً:} «والصواب من القول... أن يُقال إنها كلمة كرهها الله لهم أن يقولوها لنبيه ﷺ... نظير الذي ذُكر عن النبي ﷺ أنه قال: \begin{hadith}لا تقولوا للعنب الكرم... ولا تقولوا عبدي ولكن قولوا فتاي\end{hadith}».
فالكلمة تحتمل معنى "الرعونة" (في لغة اليهود)، وتحتمل "أرعنا سمعك" (في لغة العرب)، فنهى الله المؤمنين عنها توقيراً للنبي ﷺ، وأمرهم بلفظ لا ريب فيه وهو (انظرنا) أي: انتظرنا وتمهل علينا لنفهم عنك.

\begin{fawaid}[قاعدة في تهذيب الألفاظ وسد الذرائع]
نهى الله المؤمنين عن استخدام لفظة (راعنا) لاحتمالها معنى فاسداً يستغله اليهود للسخرية بالنبي ﷺ، رغم أن المؤمنين يقصدون بها معنى صحيحاً. وهذا أصل في وجوب اختيار أطايب الكلام، وتجنب الألفاظ الموهمة للباطل، سداً لذريعة الفساد وسخرية الأعداء.
\end{fawaid}

\separator

\section{تأويل (ما يود الذين كفروا من أهل الكتاب ولا المشركين أن ينزل عليكم من خير من ربكم)}
\isnad{قال أبو جعفر:} «ما يحب الكافرون من أهل الكتاب ولا من المشركين بالله... أن ينزل عليكم شيئاً من الخير الذي هو عنده... وانما أحبت اليهود وأتباعهم من المشركين ذلك حسداً وبغياً منهم على المؤمنين».

\begin{fawaid}[استنباط: الاطلاع على بواطن الأعداء]
بيّن الطبري أن في هذه الآية دلالة بيّنة على نهي المؤمنين عن الركون لأعدائهم والاستماع لقولهم وقبول نصيحتهم، لأن الله أطلع المؤمنين على ما يستبطنه الكفار من الضغينة والحسد، وإن أظهروا بألسنتهم خلاف ذلك، لئلا يغتروا بظاهرهم.
ومن اللطائف ما قاله \rawi{ابن عطية} أن "الخير" هنا يقصد به توقير النبي ﷺ، فإن تعظيم المسلمين لنبيهم يغيظ الكفار جداً.
\end{fawaid}

\separator

\section{تأويل (ما ننسخ من آية أو ننسها نأت بخير منها أو مثلها)}
\isnad{قال أبو جعفر:} «ما ننقل من حكم آية إلى غيره نغيره ونبدله... ولا يكون ذلك إلا في الأمر والنهي والحظر والإطلاق... فأما الأخبار فلا يكون منها ناسخ ولا منسوخ».

\isnad{وأصل النسخ:} نقل الكتاب، فنسخ الحكم نقله وتحويل العباد عنه إلى غيره.

واختلفت القراءة في قوله: \begin{ayah}أَوْ نُنسِهَا\end{ayah}؛ فقرأها عامة قراء الحجاز والعراق: (نُنسِهَا) بضم النون، بمعنى نتركها أو نُمحها من القلوب. وقرأها آخرون: (نَنْسَأْهَا) بفتح النون والهمز، بمعنى نؤخرها.
\isnad{قال الطبري مرجحاً:} «وأولى القراءات بالصواب... من قرأ (أو نُنسها) بمعنى نتركها... فإذا قُرئت كذلك بالمعنى الذي وصفت، فهو يشتمل على معنى الإنساء الذي هو بمعنى الترك... ومعنى النسء الذي هو بمعنى التأخير».

\begin{fawaid}[منهج الطبري في الترجيح وجمع المعاني]
هذا الموضع من أهم مواضع التفسير التي تبرز منهج الطبري في الجمع بين القراءات والترجيح، فهو يختار القراءة التي تستوعب دلالات المعاني المختلفة متى أمكن ذلك، ويرد القراءات الشاذة المخالفة لرسم المصحف وإجماع القرأة.
\end{fawaid}

\isnad{ورَدَّ الطبري} على من أنكر جواز نسيان النبي ﷺ لشيء من القرآن استدلالاً بـ (أو ننسها)، مبيناً أنه يجوز أن ينسي الله نبيه ما يشاء مما نسخه ورفع تلاوته، كما قال: \begin{ayah}سَنُقْرِئُكَ فَلَا تَنسَىٰ ۖ إِلَّا مَا شَاءَ اللَّهُ\end{ayah}.

\separator

\section{تأويل (أم تريدون أن تسألوا رسولكم كما سئل موسى من قبل)}
اختلفوا في سبب نزولها؛ فقيل نزلت حين سأل رافع بن حريملة وغيره النبي ﷺ أن يأتيهم بكتاب من السماء ويفجر لهم أنهاراً (كما سُئل موسى رؤية الله جهرة)، وقيل سألوا أن يجعل لهم الصفا ذهباً، وقيل سألوا كفارات ككفارات بني إسرائيل.

وناقش \rawi{الطبري} معنى (أَمْ) هنا، فرجح أنها "استفهام مبتدأ"، وليست متصلة بما قبلها للشك، وأنها تدل على توبيخهم على سؤال التعنت.

\begin{fawaid}[دلالة (أم) المنقطعة]
تأتي (أم) في لغة العرب للاستفهام المستأنف (المنقطعة) بمعنى (بل)، وتكون توبيخاً وإنكاراً، ولا يُشترط أن تكون دائماً متصلة (للتسوية) أو مبنية على الشك.
\end{fawaid}

\separator

بهذا القدر نكتفي، ونسأل الله أن ينفعنا بما قرأنا، ونكمل في الدرس القادم بإذن الله. والسلام عليكم ورحمة الله وبركاته.